% Integrated strengthening, based on the frozen proof revision of 2026-10-06.
% No preamble. All valuations have v_p(0)=+infinity.
\section{Threshold-free reciprocal transfer at prime-power section orders}
\label{sec:prime-power-transfer-final}

The completion of Nishibuchi's conjecture fixes the requested eventual
congruence. The stronger question is whether the sufficient starting row
can be removed, and how much precision the difference order and the step
actually provide. This section answers the first question at every
prime-power section order. Subsequent sections give a factorial improvement
and exact results at section order four.

Throughout these sections, the normalized source and its reciprocal are
\begin{equation}
 G_{N,j}(t)=j!\sum_{n\geq0}\frac{t^{Nn}}{(Nn+j)!},\qquad
 F_{N,j}(t)=G_{N,j}(t)^{-1}
   =\sum_{n\geq0}C_{N,j}(n)\frac{t^{Nn}}{(Nn)!},
 \label{eq:final-normalization}
\end{equation}
where $N\geq1$ and $0\leq j<N$. In particular, $C_{N,j}(0)=1$.
The coefficients in the original convention satisfy
\begin{equation}
 E_{Nn}^{(N,j)}=j!C_{N,j}(n),\qquad
 E_d^{(N,j)}=0\quad(N\nmid d).
 \label{eq:final-original-normalization}
\end{equation}
Thus every valuation of a nonzero difference gains $v_p(j!)$ on returning
to the original coefficients. This factor must be retained when stating
optimal periods or the final congruence.

Write $S f(n)=f(n+1)$ and $\Delta_H=S^H-1$. A mixed difference
$\Delta_{H_1}\cdots\Delta_{H_k}$ acts on the row variable $n$, rather than
on the full Euler subscript. We use $v_p(0)=+\infty$. Orders and starting
rows are nonnegative integers. Steps are positive integers unless stated
otherwise. An empty mixed difference is the original value, and a zero
ordinary step makes every nonempty product containing it zero.

\subsection{The Hurwitz lattice and translated differential flatness}

\begin{definition}[Hurwitz integral series]
For a prime $p$, let
\[
 \mathcal H_p=\left\{\sum_{d\geq0}f_d\frac{t^d}{d!}:
                    f_d\in\mathbb Z_p\right\}
       \subset\mathbb Q_p[[t]].
\]
For $N=p^a$ and $\lambda\in\mathbb Z_p$, put
\[
 D=\frac{d}{dt},\qquad L_\lambda=\frac{D^N-\lambda}{p}.
\]
A series $A$ is $\lambda$-flat if $L_\lambda^k A\in\mathcal H_p$ for every
$k\geq0$.
\end{definition}

Products preserve $\mathcal H_p$ because the coefficient at factorial
degree $d$ is
\[
 (AB)_d=\sum_{i=0}^d\binom di a_i b_{d-i}.
\]
The derivative $D$ acts by the coefficient shift $(DA)_d=a_{d+1}$.
Consequently every nonnegative power of $D$ preserves the lattice.
If $A\in\mathcal H_p$ has unit constant term, its reciprocal belongs to
the same lattice: with $B=A^{-1}$,
\begin{equation}
 b_0=a_0^{-1},\qquad
 b_d=-a_0^{-1}\sum_{i=1}^d\binom di a_i b_{d-i}.
 \label{eq:hurwitz-reciprocal-recurrence}
\end{equation}
If $A$ is supported at degrees divisible by $N$, this recurrence also
proves that $B$ has the same support.

\begin{lemma}[Translated product formula]
\label{lem:translated-hurwitz-product}
For $N=p^a$, $a\geq1$, and $\lambda,\mu\in\mathbb Z_p$,
\begin{equation}
 L_{\lambda+\mu}(AB)
  =(L_\lambda A)B+A(L_\mu B)
       +\sum_{i=1}^{N-1}c_i(D^iA)(D^{N-i}B),
 \qquad c_i=\frac1p\binom Ni\in\mathbb Z.
 \label{eq:translated-product}
\end{equation}
Flatness types add under multiplication.
\end{lemma}

\begin{proof}
The identity $(1+X)^{p^a}=1+X^{p^a}$ over $\mathbb F_p$ shows that every
interior binomial coefficient $\binom Ni$ is divisible by $p$.
Subtract $(\lambda+\mu)AB$ from the ordinary $N$th derivative product
formula and divide by $p$ to obtain \eqref{eq:translated-product}.

For completeness, the iterated formula is an integral multinomial
formula, rather than a product rule with unspecified remainders. On the
tensor of the two factors use the mutually commuting operators
\[
 L_\lambda\otimes1,\quad 1\otimes L_\mu,\quad
 c_iD^i\otimes D^{N-i}\quad(1\leq i<N).
\]
Applying multiplication after the $k$th power of their sum gives exactly
$L_{\lambda+\mu}^k(AB)$. Expanding it yields terms indexed by
$a,b,e_1,\ldots,e_{N-1}\geq0$, with $a+b+\sum_i e_i=k$:
\begin{equation}
 \frac{k!}{a!b!\prod_i e_i!}\prod_i c_i^{e_i}
 (D^\alpha L_\lambda^a A)(D^\beta L_\mu^b B),
 \quad
 \alpha=\sum_i i e_i,\quad
 \beta=\sum_i(N-i)e_i.
 \label{eq:translated-multinomial}
\end{equation}
Every scalar is integral. If $A$ and $B$ have the stated flatness types,
both factors in every term belong to $\mathcal H_p$, as does their product.
\end{proof}

\begin{lemma}[Reciprocal transfer]
\label{lem:translated-hurwitz-inverse}
If $A$ is $\lambda$-flat and has unit constant term, then $A^{-1}$ is
$(-\lambda)$-flat.
\end{lemma}

\begin{proof}
Set $B=A^{-1}$. Its order-zero integrality was established in
\eqref{eq:hurwitz-reciprocal-recurrence}. Assume by induction that
$L_{-\lambda}^bB\in\mathcal H_p$ for $b<k$. Since $AB=1$,
\[
 L_0^k(AB)=\left(\frac{D^N}{p}\right)^k1=0\qquad(k>0).
\]
Use \eqref{eq:translated-multinomial} with $\mu=-\lambda$.
Exactly one term has inverse order $b=k$: it has $a=0$ and all $e_i=0$,
and is $A L_{-\lambda}^k B$. In every other term $b<k$, so its second
factor is integral by induction and derivative-shift integrality. Its
first factor is integral by the flatness of $A$, and its scalar is
integral. Hence $A L_{-\lambda}^k B\in\mathcal H_p$. Multiplying by the
already integral $B$ proves $L_{-\lambda}^k B\in\mathcal H_p$.
This induction takes place in the entire Hurwitz lattice. It places no
restriction on the coefficient degree or the starting row.
\end{proof}

\subsection{The actual Euler source}

The packed factorial coefficient of $G_{N,j}$ is
\begin{equation}
 g_{N,j}(n)=\frac{j!(Nn)!}{(Nn+j)!}
   =\prod_{i=1}^j\left(1+\frac Ni n\right)^{-1}.
 \label{eq:actual-prime-power-kernel}
\end{equation}
The empty product is one. If $N=p^a$ and $1\leq i<N$, then
$v_p(i)<a$, so $N/i\in p\mathbb Z_p$. Every affine denominator in
\eqref{eq:actual-prime-power-kernel} is therefore a unit on $\mathbb Z_p$.

For $s\in p\mathbb Z_p$, an exact finite-difference identity is
\begin{equation}
 \Delta_1^k\frac1{1+sn}
  =\frac{(-s)^k k!}{\prod_{h=0}^k(1+s(n+h))}.
 \label{eq:affine-reciprocal-exact}
\end{equation}
It follows by induction, subtracting adjacent rational functions. Each
successive subtraction adds the factor $-(k+1)s$ in the numerator.
Its denominator is a unit, so it belongs to $p^k\mathbb Z_p$.
The discrete Leibniz formula
\begin{equation}
 \Delta_1^k(fg)(n)
 =\sum_{r=0}^k\binom kr
       (\Delta_1^r f)(n)(\Delta_1^{k-r}g)(n+r)
 \label{eq:discrete-leibniz-final}
\end{equation}
then proves $\Delta_1^k g_{N,j}(n)\in p^k\mathbb Z_p$ for every $k,n$.
The coefficient at degree $Nn$ of $L_1^kG_{N,j}$ is precisely
$p^{-k}\Delta_1^k g_{N,j}(n)$, and the other residue classes vanish.
Thus the actual source is $1$-flat.

\begin{theorem}[All-row prime-power transfer]
\label{thm:all-prime-power-transfer}
Let $p$ be prime, $a\geq1$, $N=p^a$, and $0\leq j<N$.
The coefficients $C_{N,j}(n)$ belong to $\mathbb Z_p$ for all $n\geq0$.
Set $A_{N,j}(n)=(-1)^n C_{N,j}(n)$ and
$T_H=S^H-(-1)^H$. For $k\geq1$ and positive $H_1,\ldots,H_k$,
\begin{align}
 v_p(\Delta_{H_1}\cdots\Delta_{H_k}A_{N,j}(n))
  &\geq k+\sum_{\nu=1}^k v_p(H_\nu),
 \label{eq:prime-power-signed-bound}\\
 v_p(T_{H_1}\cdots T_{H_k}E_{Nn}^{(N,j)})
  &\geq v_p(j!)+k+\sum_{\nu=1}^k v_p(H_\nu).
 \label{eq:prime-power-original-twist}
\end{align}
At $p=2$, the ordinary differences of $C_{N,j}$ and of the original
coefficients satisfy the same bounds. These dyadic statements also hold
for $N=1$.
\end{theorem}

\begin{proof}
Reciprocal transfer gives $L_{-1}^kF_{N,j}\in\mathcal H_p$.
Extraction of the packed coefficient gives
$(S+1)^k C_{N,j}(n)\in p^k\mathbb Z_p$. Moreover,
\[
 \Delta_1^k A_{N,j}(n)=(-1)^{n+k}(S+1)^k C_{N,j}(n),
\]
so $A_{N,j}$ has the ordinary $p^k$ difference property.
At $p=2$, $L_1=L_{-1}-1$. Its integral binomial expansion on
$F_{N,j}$ gives $\Delta_1^k C_{N,j}(n)\in2^k\mathbb Z_2$ directly.
For $N=1$ the interior sum in the product formula is empty, so the same
argument applies without change. Equivalently, $F_{1,0}(t)=e^{-t}$.

To retain the valuation of a step, let $B$ be any sequence satisfying
$\Delta_1^m B\in p^m\mathbb Z_p$ at every order and row. Use the finite
operator identity
\[
 \Delta_H=\sum_{i=1}^H\binom Hi\Delta_1^i.
\]
The relation $i\binom Hi=H\binom{H-1}{i-1}$ implies
\[
 v_p\left(p^i\binom Hi\right)
 \geq i+v_p(H)-v_p(i)\geq1+v_p(H).
\]
Multiplying these identities for the $k$ steps, every term of total
order $M$ has its scalar multiplied by $p^M$, and is divisible by
$p^{k+\sum_\nu v_p(H_\nu)}$. This proves the mixed-step bound.
Finally,
\[
 \Delta_H A_{N,j}(n)=(-1)^{n+H}T_H C_{N,j}(n).
\]
Iteration contributes only a global sign. Multiplying by $j!$ proves
\eqref{eq:prime-power-original-twist}.
\end{proof}

\begin{corollary}[Threshold-free dyadic congruence]
\label{cor:all-dyadic-final}
For $N=2^a$, $a\geq0$, all permitted $j$, and $n\geq0$,
\begin{equation}
 v_2(\Delta_{H_1}\cdots\Delta_{H_k}E_{Nn}^{(N,j)})
  \geq v_2(j!)+k+\sum_{\nu=1}^k v_2(H_\nu).
 \label{eq:all-dyadic-final-bound}
\end{equation}
In particular, $\Delta_H^k$ gains $k(1+v_2(H))$ bits before normalization.
For $r>v_2(j!)$, the row displacement
$H=2^{\max(0,r-v_2(j!)-1)}$ is a sufficient period modulo $2^r$ from
row zero. If $r\leq v_2(j!)$, all original coefficients vanish modulo
$2^r$.
\end{corollary}

\begin{remark}[Necessary boundaries]
At an odd prime the ordinary step-one difference cannot replace the
sign-twisted operator in the general theorem. For $j=0$,
$C_{N,0}(0)=1$ and $C_{N,0}(1)=-1$, so their ordinary difference is
$-2$, a unit at every odd prime. The proof also does not extend
unchanged to non-prime-power $N$, whose interior binomial coefficients
need not be divisible by the chosen prime. Neither observation is a
counterexample to Nishibuchi's original eventual congruence.
\end{remark}

\section{Factorial precision at section order four}
\label{sec:factorial-four-final}

Write $C_j=C_{4,j}$ and $F_j=F_{4,j}$. The all-row theorem removes the
threshold, but at $N=4$ it leaves a further factorial gain unused.
The exact algebra below proves that gain for all four addresses and
identifies several infinite families where it is sharp.

\begin{theorem}[Four-address factorial bound]
\label{thm:four-factorial-bound}
For $j\in\{0,1,2,3\}$, $n\geq0$, $k\geq1$, and positive mixed steps,
\begin{equation}
 v_2(\Delta_{H_1}\cdots\Delta_{H_k}C_j(n))
 \geq k+v_2(k!)+\sum_{\nu=1}^k v_2(H_\nu)+\beta_j,
 \quad (\beta_0,\beta_1,\beta_2,\beta_3)=(0,0,3,1).
 \label{eq:factorial-four-lower}
\end{equation}
Address two also satisfies
\begin{equation}
 v_2(\Delta_{H_1}\cdots\Delta_{H_k}C_2(n))
 \geq 2k+v_2(k!)+\left\lfloor\frac{k}{2}\right\rfloor
             +\sum_{\nu=1}^k v_2(H_\nu).
 \label{eq:factorial-four-j2-filtration}
\end{equation}
For the original coefficients, the offsets in
\eqref{eq:factorial-four-lower} are $(0,0,4,2)$.
No positive-order offset is asserted for the empty difference.
\end{theorem}

The proof of the first two addresses uses one classical arithmetic
input: Frobenius's exact congruence for the ordinary Euler numbers in
the $\operatorname{sech}$ convention, as stated and derived by
Gessel~\cite{gessel}, Section~9. The address-two and address-three
bounds below are finite Bernoulli calculations and do not require that
input. A divided-power argument is necessary: ordinary integral moments
alone do not justify division by $k!$.

\subsection{Euler moments and an integral divided-power algebra}

Let $\mathcal L$ be the ordinary Euler evaluation functional on a
variable $\mathcal E$, defined by
\[
 \mathcal L(\mathcal E^d)=\operatorname{Euler}_d,
 \qquad \sum_{d\geq0}\operatorname{Euler}_d\frac{t^d}{d!}
                  =\operatorname{sech}t.
\]
The functional is linear. It is not a multiplicative evaluation.
Independent variables mean that moments of polynomials in separate
variables factor. Odd moments vanish. Frobenius's congruence is
\begin{equation}
 v_2\left(\sum_{r=0}^m(-1)^{m-r}\binom mr
                 \operatorname{Euler}_{2n+2r}\right)
       =m+v_2(m!)\qquad(n,m\geq0).
 \label{eq:classical-frobenius-exact}
\end{equation}
For $m=0$, the assertion is that every even Euler number is odd.
With $P=(\mathcal E^2-1)/2$ and
$\gamma_m(P)=P^m/m!$, the case $n=0$ says
\begin{equation}
 \mathcal L(\gamma_m(P))\in\mathbb Z_2^{\times}
                 \qquad(m\geq0).
 \label{eq:euler-pd-odd-moments}
\end{equation}
Introduce an independent Euler variable $\mathcal F$ and
$Q=(\mathcal F^2-1)/2$.

Put $K=\mathbb Q_2(i)$, $R=\mathbb Z_2[i]$,
$a=(1+i)/2$, and $b=(1-i)/2$. The exact identity
\begin{equation}
 F_0(t)=\frac{2}{\cosh t+\cos t}
       =\operatorname{sech}(at)\operatorname{sech}(bt)
 \label{eq:four-euler-factorization}
\end{equation}
gives, for $U=a\mathcal E+b\mathcal F$,
$\mathcal L(U^{4n})=C_0(n)$. Although $a$ and $b$ have negative
valuation, their denominators disappear from the following exact
polynomial identities:
\begin{align}
 U^2&=\mathcal E\mathcal F+i(P-Q),\nonumber\\
 H_0:=\frac{U^4-1}{2}
   &=P+Q-\frac{P^2}{2}-\frac{Q^2}{2}
                 +3PQ+i\mathcal E\mathcal F(P-Q).
 \label{eq:four-euler-H}
\end{align}

Let $\Gamma_R(P,Q)$ be the $R$-subalgebra generated by all
$\gamma_m(P)$ and $\gamma_r(Q)$. Adjoin $\mathcal E,\mathcal F$ with
\[
 \mathcal E^2=1+2P,\qquad \mathcal F^2=1+2Q,
\]
and call the resulting algebra $\mathcal A_0$. Its elements are finite
expressions. No completion is being presumed. It is generated as a
$\Gamma_R(P,Q)$-module by $1,\mathcal E,\mathcal F,\mathcal E\mathcal F$.
The functional maps this module to $R$. For an even basis term this is
\eqref{eq:euler-pd-odd-moments} and independence. For a basis term with
an odd Euler factor the moment is zero, since $P$ and $Q$ are even
polynomials in their respective variables.

Every displayed component of $H_0$ has all its divided powers in
$\mathcal A_0$. Besides the generators, the needed identities are
\begin{align}
 \gamma_m(P^2/2)
     &=\frac{(2m)!}{2^m m!}\gamma_{2m}(P),\nonumber\\
 \gamma_m(3PQ)&=3^m m!\gamma_m(P)\gamma_m(Q),\nonumber\\
 \gamma_m(P-Q)&=\sum_{r=0}^m(-1)^r
                   \gamma_{m-r}(P)\gamma_r(Q).
 \label{eq:components-divided-powers}
\end{align}
The first scalar counts partitions into pairs and is an integer.
For the last term of $H_0$, multiply the third identity by
$(i\mathcal E\mathcal F)^m$ and reduce the Euler powers using the
integral quadratic relations. Finally,
\[
 \gamma_k(X+Y)=\sum_{r=0}^k\gamma_r(X)\gamma_{k-r}(Y)
\]
proves $\gamma_k(H_0)\in\mathcal A_0$. Applying the functional and
putting $W=U^4-1=2H_0$ gives
\begin{equation}
 \mathcal L(W^m)\in2^m m!R\qquad(m\geq0).
 \label{eq:euler-W-factorial}
\end{equation}
These moments are rational. The identity
$R\cap\mathbb Q_2=\mathbb Z_2$, obtained from the linear independence
of $1,i$ over $\mathbb Q_2$, descends every such divisibility statement
to $\mathbb Z_2$. No field property of $R/2R$ is needed.

For every starting row,
\begin{equation}
 \Delta_1^k C_0(n)
   =\mathcal L(U^{4n}W^k)
   =\sum_{s=0}^n\binom ns\mathcal L(W^{k+s}).
 \label{eq:base-row-finite-expansion}
\end{equation}
Each summand lies in $2^k k!\mathbb Z_2$, proving the factorial
bound for address zero from row zero onward.

\subsection{A finite step-expansion lemma}

The conversion from unit steps to arbitrary mixed steps is used several
times, so it is useful to make it explicit. Write $\Delta=\Delta_1$.
For $H>0$,
\begin{equation}
 \Delta_H=\sum_{m=1}^H\binom Hm\Delta^m.
 \label{eq:finite-step-expansion}
\end{equation}
A mixed term with local orders $m_1,\ldots,m_k\geq1$ and
$M=\sum_\nu m_\nu$ therefore has valuation at least
\[
 M+v_2(M!)+\beta
    +\sum_\nu\bigl(v_2(H_\nu)-v_2(m_\nu)\bigr)
\]
if the unit-order estimate has offset $\beta$ for every positive order.
The inequalities
\[
 m-1\geq v_2(m),\qquad M\geq k,\qquad v_2(M!)\geq v_2(k!)
\]
make this at least $k+v_2(k!)+\beta+\sum_\nu v_2(H_\nu)$.
Every expansion is finite. This proves the desired mixed-step transfer
without an exchange of limits.

\subsection{An elementary Bernoulli denominator bound}

Let $B_m$ denote the Bernoulli numbers with $B_1=-1/2$.
The identity obtained by substituting $u=e^t-1$ into
$\log(1+u)/u$ is
\begin{equation}
 B_m=\sum_{r=0}^m(-1)^r\frac{r!}{r+1}
                                  \left\{\begin{matrix}m\\r\end{matrix}\right\},
 \label{eq:bernoulli-stirling-denominator}
\end{equation}
where the braces are Stirling numbers of the second kind. Each
$2r!/(r+1)$ is dyadically integral. Indeed, if $s=v_2(r+1)>0$, then
$r\geq2^s-1$ and
\[
 v_2(r!)\geq v_2((2^s-1)!)=2^s-1-s\geq s-1.
\]
The last inequality is $2^s\geq2s$. The case $s=0$ is immediate.
Consequently
\begin{equation}
 2B_m\in\mathbb Z_2\qquad(m\geq0).
 \label{eq:elementary-bernoulli-bound}
\end{equation}
Thus the transfer factors below require no appeal to the
von Staudt--Clausen theorem.

\subsection{Address one and the scaled-binomial transfer}

The address-one source is
$G_{4,1}(t)=(\sinh t+\sin t)/(2t)$. The identities $a+b=1$,
$a-b=i$, and $ab=1/2$ give
\begin{equation}
 D_1(t):=G_{4,1}(t)F_0(t)
 =\frac12\left(\frac{\tanh(at)}{at}
                         +\frac{\tanh(bt)}{bt}\right).
 \label{eq:address-one-tangent-factor}
\end{equation}
For example, its numerator follows from
\[
 \sinh t+\sin t
 =2b\sinh(at)\cosh(bt)+2a\cosh(at)\sinh(bt),
\]
while $\cosh t+\cos t=2\cosh(at)\cosh(bt)$.
Only degrees divisible by four survive. Its actual Hurwitz packets are
\begin{equation}
 D_1(t)=\sum_{i\geq0}d_i\frac{t^{4i}}{(4i)!},\qquad
 d_i=\left(-\frac14\right)^i
 \frac{2^{4i+2}(2^{4i+2}-1)B_{4i+2}}{(4i+2)(4i+1)}.
 \label{eq:address-one-packets}
\end{equation}
The constant is one. Since $v_2(4i+2)=1$, the elementary Bernoulli bound
shows $v_2(d_i)\geq2i$. Hurwitz inversion therefore gives
\[
 D_1(t)^{-1}=\sum_{i\geq0}h_i\frac{t^{4i}}{(4i)!},
 \qquad h_0=1,\quad v_2(h_i)\geq2i.
\]
This follows by induction from the inverse recurrence: for packet index
$i$, every product $d_rh_{i-r}$ has valuation at least $2i$, and its
binomial multiplier is integral.
Because $F_1=F_0D_1^{-1}$,
\begin{equation}
 C_1(n)-C_0(n)
     =\sum_{i\geq1}h_i\binom{4n}{4i}C_0(n-i).
 \label{eq:address-one-exact-transfer}
\end{equation}
At a natural $n$ this is a finite sum, with $i\leq n$.

Here is the polynomial estimate that supplies one additional bit.
For $i\geq1$, put $P_i(x)=4^i\binom{4x}{4i}$. Its constant is zero,
and its degree-$d$ coefficient is
\begin{equation}
 [x^d]P_i(x)=\frac{4^i4^d s(4i,d)}{(4i)!},
 \label{eq:scaled-binomial-coefficients}
\end{equation}
where $s(m,d)$ is a signed Stirling number of the first kind.
In each product contributing to $s(4i,d)$, at most $d-1$ of the
$2i-1$ even integers in $1,\ldots,4i-1$ can be omitted. Thus
\[
 v_2(s(4i,d))\geq\max(0,2i-d).
\]
Writing $s_2(i)$ for the number of ones in the binary expansion of $i$,
and using $v_2((4i)!)=4i-s_2(i)$, gives
\begin{equation}
 v_2([x^d]P_i)
 \geq2d-2i+s_2(i)+\max(0,2i-d)
 \geq d+s_2(i)\geq d+1\quad(d\geq1).
 \label{eq:scaled-binomial-degree-bound}
\end{equation}
For $k\geq1$, $\Delta_1^k x^d$ vanishes when $d<k$ and is an
integer polynomial with values divisible by $k!$. To see the latter,
expand $x^d=\sum_r\left\{\begin{smallmatrix}d\\r\end{smallmatrix}\right\}
r!\binom xr$ and apply
$\Delta_1^k\binom xr=\binom{x}{r-k}$.
It follows that
\begin{equation}
 \Delta_1^kP_i(x)\in2^{k+1}k!\mathbb Z_2
                      \qquad(k\geq0,\ x\in\mathbb Z_2).
 \label{eq:scaled-binomial-difference-bound}
\end{equation}
For $k=0$ the bound follows directly from the coefficient estimate.

To apply discrete products in \eqref{eq:address-one-exact-transfer}
without a moving truncation, first extend $C_0$ to $\mathbb Z_2$ by its
uniformly convergent Newton series. The all-row bound
$v_2(\Delta_1^m C_0(0))\geq m$ proves this convergence. Its factorial
bounds extend to every dyadic starting point by continuity and
approximation with natural rows. The sum in
\eqref{eq:address-one-exact-transfer} is then uniformly convergent on
$\mathbb Z_2$: both the binomial polynomial and $C_0(x-i)$ are
integral-valued and $v_2(h_i)\geq2i$.

For the pointwise product $P_i(x)C_0(x-i)$, the discrete Leibniz formula
has scalar divisibility
\[
 \binom kr(2^{r+1}r!)(2^{k-r}(k-r)!)=2^{k+1}k!.
\]
Multiplication by $h_i/4^i\in\mathbb Z_2$ and uniform summation now give
\begin{equation}
 v_2(\Delta_1^k(C_1-C_0)(n))
               \geq k+v_2(k!)+1\qquad(k,n\geq0).
 \label{eq:address-one-extra-bit}
\end{equation}
The same extra bit holds for positive mixed steps by
\eqref{eq:finite-step-expansion}. Combining this with the address-zero
bound proves Theorem~\ref{thm:four-factorial-bound} for address one.

\subsection{Central factorials and integral Bernoulli divided powers}

\begin{lemma}[Bernoulli divided-power moments]
\label{lem:bernoulli-central-pd}
Let $Z$ be the ordinary Bernoulli variable,
$\mathcal L(Z^r)=B_r$, with exponential transform $t/(e^t-1)$.
For $P=Z(Z+1)$ and $\mu_m=\mathcal L(P^m/m!)$,
\begin{equation}
 \mu_m\in\mathbb Z_2,\qquad
 v_2(\mu_m)\geq s_2(m)-1\quad(m\geq1).
 \label{eq:bernoulli-central-moment-bound}
\end{equation}
Moreover, $\mu_m$ is odd exactly when $m=0$ or $m$ is a positive power
of two.
\end{lemma}

\begin{proof}
Define the central-factorial basis polynomials
\[
 P_r(X)=\prod_{a=0}^{r-1}(X-a(a+1)).
\]
At $X=Z(Z+1)$, they have the exact factorization
\begin{equation}
 P_r(Z(Z+1))
   =\prod_{a=0}^{r-1}(Z-a)(Z+a+1)
   =(2r)!\binom{Z+r}{2r}.
 \label{eq:central-factorial-substitution}
\end{equation}
The Bernoulli transform gives
$\mathcal L((1+T)^Z)=\log(1+T)/T$, whence
$\mathcal L(\binom Zj)=(-1)^j/(j+1)$. Finite Vandermonde expansion
therefore gives
\begin{align}
 \mathcal L\left(\binom{Z+r}{2r}\right)
 &=(-1)^r\sum_{q=0}^r(-1)^q\binom rq\frac1{r+q+1}\nonumber\\
 &=(-1)^r\frac{(r!)^2}{(2r+1)!}.
 \label{eq:central-factorial-evaluation}
\end{align}
The last equality is the rational identity obtained by expanding
$\int_0^1t^r(1-t)^r\,dt$. This use of an ordinary finite integral makes
no assertion about a positive or bounded $p$-adic measure. In
particular,
\[
 \mathcal L(P_r(P))=(-1)^r\frac{(r!)^2}{2r+1}.
\]

Expand $X^m=\sum_{r=0}^mT(m,r)P_r(X)$. Since
$XP_r(X)=P_{r+1}(X)+r(r+1)P_r(X)$, these integer coefficients satisfy
\begin{equation}
 T(m+1,r)=T(m,r-1)+r(r+1)T(m,r),
 \quad T(0,0)=1.
 \label{eq:central-factorial-recurrence}
\end{equation}
Every $r(r+1)$ is even. Induction, including the boundary indices,
proves $2^{m-r}\mid T(m,r)$ for $0\leq r\leq m$.
For $m>0$ the $r=0$ term vanishes. Each remaining summand of
$\mathcal L(P^m/m!)$ has valuation at least
\begin{equation}
 m-r+2v_2(r!)-v_2(m!)
   =r+s_2(m)-2s_2(r)\geq s_2(m)-1\geq0.
 \label{eq:central-factorial-term-bound}
\end{equation}
Here $r\geq2s_2(r)-1$ follows from
$r\geq2^{s_2(r)}-1$ and $2^u\geq2u$ for $u\geq1$.
This proves integrality and the stronger bound.

If $m$ is not a power of two, $s_2(m)\geq2$ makes $\mu_m$ even.
If $m$ is a positive power of two, the only central indices which can
contribute an odd residue are $r=1$ and $r=3$, since
$r-2s_2(r)=-1$ only at these two positive integers. The $r=1$ term is
odd: $T(m,1)=2^{m-1}$ and $v_2(m!)=m-1$. For $m\geq4$ a power of
two, the $r=3$ term is even. More explicitly,
\[
 \frac{T(m,3)}{2^{m-3}}=h_{m-3}(1,3,6),
 \qquad
 \sum_{d\geq0}h_d(1,3,6)t^d
       =\frac1{(1-t)(1-3t)(1-6t)}.
\]
Reduction modulo two gives $(1-t)^{-2}$, whose degree $m-3$
coefficient has parity $m-2=0$. At $m=1,2$ the $r=3$ term is absent,
and the $r=1$ term is again odd. Finally $\mu_0=1$.
\end{proof}

\subsection{Address two and its stronger filtration}

The variable $\mathcal B=2Z+1$ has transform
\[
 e^t\frac{2t}{e^{2t}-1}=\frac{t}{\sinh t}.
\]
Its odd moments vanish, and
$P=(\mathcal B^2-1)/4=Z(Z+1)$ has the divided-power moments of
Lemma~\ref{lem:bernoulli-central-pd}. Introduce an independent copy
$\mathcal D$ and $Q=(\mathcal D^2-1)/4$. The exact factorization
\begin{equation}
 F_2(t)=\frac{t^2}{\cosh t-\cos t}
       =\frac{at}{\sinh(at)}\frac{bt}{\sinh(bt)}
 \label{eq:address-two-factorization}
\end{equation}
shows that $U=a\mathcal B+b\mathcal D$ satisfies
$\mathcal L(U^{4n})=C_2(n)$. Direct expansion gives
\begin{equation}
 H_2:=\frac{U^4-1}{4}
     =P+Q-P^2-Q^2+6PQ+i\mathcal B\mathcal D(P-Q).
 \label{eq:address-two-H}
\end{equation}

Let $\mathcal A_2$ be the divided-power algebra over $R$ generated by
$P,Q$, adjoining $\mathcal B,\mathcal D$ with
$\mathcal B^2=1+4P$, $\mathcal D^2=1+4Q$.
The functional is integral on its module basis
$1,\mathcal B,\mathcal D,\mathcal B\mathcal D$ by the Bernoulli lemma,
independence, and vanishing of odd moments. Every component in
\eqref{eq:address-two-H} has integral divided powers, by the same
finite identities as before. $P^2$ is even easier than $P^2/2$.
Consequently
\begin{equation}
 \mathcal L((U^4-1)^m)\in4^m m!\mathbb Z_2,
 \qquad
 v_2(\Delta_1^kC_2(n))\geq2k+v_2(k!).
 \label{eq:address-two-basic-factorial}
\end{equation}
The starting-row conclusion follows from the finite base expansion
\eqref{eq:base-row-finite-expansion}, now with $W=4H_2$.

The low Newton coefficients of the actual reciprocal are
\begin{equation}
 C_2(0)=1,\quad C_2(1)=-\frac1{15},\quad C_2(2)=\frac{13}{45},
 \qquad a_1=-\frac{16}{15},\quad a_2=\frac{64}{45},
 \label{eq:address-two-initial-newton}
\end{equation}
where $a_m=\Delta_1^m C_2(0)$. Thus $v_2(a_1)=4$ and $v_2(a_2)=6$.
For $m\geq3$, \eqref{eq:address-two-basic-factorial} bounds $a_m$ by
$2m+v_2(m!)$. The finite shifted Newton formula
\begin{equation}
 \Delta_1^kC_2(n)=\sum_{m=k}^{n+k}a_m\binom n{m-k}
 \label{eq:shifted-newton-finite}
\end{equation}
therefore gives
$v_2(\Delta_1^kC_2(n))\geq k+v_2(k!)+3$ for every $k\geq1$.
For $k\geq3$ this is $2k\geq k+3$. For $k=1,2$ it uses the exact
initial coefficients and the higher tail bound. The finite step
expansion proves the offset three in
\eqref{eq:factorial-four-lower}.

There is a still stronger filtration. Put
\[
 X=P+Q,\qquad \tau=i\mathcal B\mathcal D,\qquad \epsilon=1+\tau.
\]
Then the exact identities are
\begin{equation}
 H_2=\epsilon X+2G,\quad
 G=-\gamma_2(P)-\gamma_2(Q)+3PQ-\tau Q,\quad
 \epsilon^2=2V,\quad V=\tau-2X-8PQ.
 \label{eq:address-two-filtration-algebra}
\end{equation}
Every divided power of $G$ is integral component by component.
The finite addition identity gives
\begin{equation}
 \gamma_k(H_2)=\sum_{r=0}^k2^r\gamma_r(G)
                  \epsilon^{k-r}\gamma_{k-r}(X)
          \in2^{\lfloor k/2\rfloor}\mathcal A_2,
 \label{eq:address-two-filtration-pd}
\end{equation}
because the $r$th term belongs to
$2^{r+\lfloor(k-r)/2\rfloor}\mathcal A_2$.
Applying the functional, rational descent, and the base-index expansion
proves the unit-step version of
\eqref{eq:factorial-four-j2-filtration}.
For a mixed term of total order $M$, its lower bound exceeds the target
at order $k$ by at least
\[
 2(M-k)+\left\lfloor\frac M2\right\rfloor
          -\left\lfloor\frac k2\right\rfloor
       +v_2(M!)-v_2(k!)-\sum_\nu v_2(m_\nu)\geq0.
\]
Indeed $\sum_\nu v_2(m_\nu)\leq M-k$. This proves the stronger mixed
bound. The earlier offset-three bound is retained, since it is stronger
at the lowest orders.

\subsection{Address three and cancellation of the apparent pole}

The address-three source is
$G_{4,3}(t)=3(\sinh t-\sin t)/t^3$. Addition identities give
\[
 \sinh t-\sin t
 =2a\sinh(at)\cosh(bt)+2b\cosh(at)\sinh(bt),
 \qquad \cosh t-\cos t=2\sinh(at)\sinh(bt).
\]
Consequently
\begin{equation}
 D_3(t):=G_{4,3}(t)F_2(t)
 =\frac32\left(\frac{\coth(at)}{at}
                            +\frac{\coth(bt)}{bt}\right).
 \label{eq:address-three-cotangent-factor}
\end{equation}
This expression is a power series: its two polar terms cancel because
$a^{-2}+b^{-2}=0$. Its constant is one, only degrees divisible by four
survive, and the Bernoulli expansion gives
\begin{equation}
 D_3(t)=\sum_{i\geq0}d_i\frac{t^{4i}}{(4i)!},\qquad
 d_i=3\left(-\frac14\right)^i
          \frac{2^{4i+2}B_{4i+2}}{(4i+2)(4i+1)}.
 \label{eq:address-three-packets}
\end{equation}
Again $v_2(d_i)\geq2i$, and the inverse coefficients $h_i$ have the same
bound. Since $F_3=F_2D_3^{-1}$,
\begin{equation}
 C_3(n)-C_2(n)=\sum_{i\geq1}h_i\binom{4n}{4i}C_2(n-i).
 \label{eq:address-three-exact-transfer}
\end{equation}
The scaled-binomial argument
\eqref{eq:scaled-binomial-difference-bound} applies unchanged.
The basic bound for $C_2$ includes order zero. Its continuous extension
is available from the threshold-free theorem. Thus the shifted sums
converge uniformly on $\mathbb Z_2$, and the pointwise product rule
proves
\begin{equation}
 v_2(\Delta_1^k(C_3-C_2)(n))\geq k+v_2(k!)+1
                                      \qquad(k,n\geq0).
 \label{eq:address-three-extra-bit}
\end{equation}
At every positive order $C_2$ has three extra bits, so this gives the
one-bit offset for $C_3$. Finite step expansion finishes the proof of
Theorem~\ref{thm:four-factorial-bound}.

\section{Sharp higher-difference families}
\label{sec:sharp-higher-final}

The lower bounds do not mean that every mixed stencil has exactly the
stated valuation. Addresses zero and one have a complete classification
of when the leading factorial precision is attained. Address two has
exact mixed precision at every positive power-of-two order.

\begin{theorem}[Exact factorial precision at addresses zero and one]
\label{thm:zero-one-exact-mixed}
Let $j\in\{0,1\}$, $k\geq1$, $n\geq0$, and $H_1,\ldots,H_k>0$.
Put
\[
 L=k+v_2(k!)+\sum_{\nu=1}^k v_2(H_\nu),\qquad
 e=\#\{\nu:2\mid H_\nu\}.
\]
The mixed difference has valuation exactly $L$ if $k$ is odd or $e$ is
even. If $k$ is even and $e$ is odd, it has valuation at least $L+1$.
In particular, for every homogeneous positive step,
\begin{equation}
 v_2(\Delta_H^kC_j(n))=k+v_2(k!)+k v_2(H).
 \label{eq:zero-one-homogeneous-exact}
\end{equation}
\end{theorem}

\begin{proof}
Return to the Euler variables of address zero and put
$X=P+Q$, $\tau=i\mathcal E\mathcal F$, $\epsilon=1+\tau$.
Direct rearrangement of \eqref{eq:four-euler-H} gives
\begin{equation}
 H_0=\gamma_2(X)+\epsilon X+2G,\qquad
 G=-\gamma_2(P)-\gamma_2(Q)+PQ-\tau Q,
 \qquad \epsilon^2=2(\tau-P-Q-2PQ).
 \label{eq:zero-parity-decomposition}
\end{equation}
The divided powers of $G$ are integral by the same component formulas
used above. Thus modulo $2\mathcal A_0$, $\gamma_k(H_0)$ equals the
divided power of $\gamma_2(X)+\epsilon X$. The identity
\[
 \gamma_r(\gamma_2(X))
   =\frac{(2r)!}{2^r r!}\gamma_{2r}(X)
\]
has an odd integer scalar. Every term involving $\epsilon^j$ for
$j\geq2$ vanishes modulo two. The surviving zero and first powers give
\begin{equation}
 \gamma_k(H_0)\equiv
       \gamma_{2k}(X)+\epsilon\gamma_{2k-1}(X)
                                  \pmod{2\mathcal A_0}.
 \label{eq:zero-parity-leading-pd}
\end{equation}
The multiplication coefficient in the second term is $2k-1$, also odd.
By the odd moments \eqref{eq:euler-pd-odd-moments} and independence,
\[
 \mathcal L(\gamma_m(P+Q))\equiv m+1\pmod2.
\]
The $\tau$-weighted moment vanishes because it contains an odd power of
each Euler variable. Therefore the first term of
\eqref{eq:zero-parity-leading-pd} evaluates to one modulo two and the
second to zero. Hence $\mathcal L(\gamma_k(H_0))$ is odd at every order,
including order zero. This proves
\[
 v_2(\Delta_1^k C_0(0))=k+v_2(k!).
\]
In the base-row expansion \eqref{eq:base-row-finite-expansion}, every
$s>0$ term has at least one extra bit after division by $2^k k!$.
Thus the same exact valuation holds at every starting row.
The extra-bit transfer \eqref{eq:address-one-extra-bit} proves it for
address one as well.

It remains to track possible ties in the mixed-step expansion.
For local orders $m_\nu\geq1$ and total $M$, the excess over $L$ is
at least
\begin{equation}
 \left(M-k-\sum_\nu v_2(m_\nu)\right)
                       +\left(v_2(M!)-v_2(k!)\right).
 \label{eq:zero-one-mixed-excess}
\end{equation}
The first bracket is nonnegative, and equality requires every
$m_\nu$ to be one or two. Two or more occurrences of two make the
second bracket positive, since the intervening factorial interval then
contains an even number. Exactly one occurrence of two can tie only
when $k$ is even. Its binomial coefficient can have the minimally
allowed valuation only if its associated step is even.

The all-one multi-index contributes one after normalization by $2^L$.
When $k$ is even, each even step supplies exactly one further tied term,
again with normalized residue one. For example,
$\binom H2=H(H-1)/2$ attains $v_2(H)-1$ exactly when $H$ is even.
All remaining factors and the exact normalized order-$k+1$ difference
are odd. Thus the total normalized parity is one for odd $k$, and
$1+e$ for even $k$. This proves the classification.
In the homogeneous case, $e$ is zero or $k$, so it is even whenever
$k$ is even, proving \eqref{eq:zero-one-homogeneous-exact}.
\end{proof}

\begin{remark}[A mixed cancellation]
The exceptional class can gain more than one bit. The exact values
\[
 \Delta_1\Delta_2C_0(0)=-33728,\qquad
 \Delta_1\Delta_2C_1(0)=-\frac{150208}{585}
\]
both have valuation six, whereas their baseline $L$ is four.
The theorem classifies attainment of the baseline. It does not give a
complete deeper valuation in the cancellation class.
\end{remark}

\begin{theorem}[Exact address-two power orders]
\label{thm:two-exact-power-orders}
For every starting row $n\geq0$ and all positive mixed steps,
\begin{equation}
 v_2(\Delta_{H_1}\cdots\Delta_{H_k}C_2(n))
       =V_2(k)+\sum_{\nu=1}^k v_2(H_\nu)
 \label{eq:two-power-orders-exact}
\end{equation}
whenever $k$ is a positive power of two, where
\[
 V_2(1)=4,\qquad V_2(2)=6,\qquad
 V_2(k)=2k+v_2(k!)+\frac{k}{2}\quad(k\geq4).
\]
For the original coefficients, the right side gains one bit.
\end{theorem}

\begin{proof}
For $k=2^r\geq4$, put $q=k/2$, which is even. Normalize the exact
expansion \eqref{eq:address-two-filtration-pd} by $2^q$ before reducing
modulo two. Its $u$th auxiliary term has excess
$u+\lfloor(k-u)/2\rfloor-q=\lfloor u/2\rfloor$.
Only auxiliary indices zero and one can therefore survive:
\begin{equation}
 \frac{\gamma_k(H_2)}{2^q}
 \equiv V^q\gamma_k(X)
       +\epsilon V^{q-1}G\gamma_{k-1}(X)
                                \pmod{2\mathcal A_2}.
 \label{eq:two-power-leading-normalized}
\end{equation}
Here $V\equiv\tau\pmod2$ and $\tau^2\equiv1\pmod2$.
Since $q$ is even and $\epsilon\tau\equiv\epsilon\pmod2$, this reduces
to $\gamma_k(X)+\epsilon G\gamma_{k-1}(X)$.
Moreover, $G\equiv\gamma_2(X)+\tau Q\pmod2$.
After evaluation, every term with a single $\tau$ vanishes, so the
second contribution reduces to
\[
 \mathcal L\bigl((\gamma_2(X)+Q)\gamma_{k-1}(X)\bigr).
\]
The $\gamma_2(X)$ product has scalar $\binom{k+1}{2}$, which is even
because $4\mid k$. The other contribution expands as
\begin{equation}
 \mathcal L(Q\gamma_{k-1}(P+Q))
    =\sum_{a=0}^{k-1}(k-a)\mu_a\mu_{k-a}.
 \label{eq:two-power-Q-evaluation}
\end{equation}
An odd summand would require $k-a$ to be odd and $\mu_{k-a}$ to be
odd. By Lemma~\ref{lem:bernoulli-central-pd}, this requires $k-a=1$.
Then $a=k-1$ is not a power of two and $\mu_a$ is even. Hence the
second contribution is zero modulo two.

For the first contribution, independence gives
$\mathcal L(\gamma_k(P+Q))=\sum_{a=0}^k\mu_a\mu_{k-a}$.
The parity generating series of these convolution moments is
\[
 (1+t+t^2+t^4+t^8+\cdots)^2
        =1+t^2+t^4+t^8+\cdots\quad\text{in }\mathbb F_2[[t]].
\]
Its coefficient is one at every $k=2^r\geq4$. Thus the normalized
moment in \eqref{eq:two-power-leading-normalized} is odd and
\[
 v_2(\Delta_1^kC_2(0))=2k+v_2(k!)+k/2.
\]
Inserting $U^{4n}=(1+4H_2)^n$ introduces only corrections with strictly
higher valuation, by the filtered bound at orders $k+s$ for $s>0$.
The equality consequently holds at every starting row.

For mixed steps the all-one local multi-index gives the claimed leading
term. A higher term of total order $M>k$ has excess at least
\[
 2(M-k)+\left\lfloor\frac M2\right\rfloor-\frac k2
       +v_2(M!)-v_2(k!)-\sum_\nu v_2(m_\nu)
       \geq M-k\geq1.
\]
Hence the leading term cannot cancel.

At order one, $a_1$ in \eqref{eq:address-two-initial-newton} has
valuation four and the Newton tail starts at valuation six. Its shifted
value remains exact. In the step expansion a local order two can lose
at most one bit through its binomial coefficient, so every correction
is still above $4+v_2(H_1)$. At order two, $a_2$ has valuation six and
the next coefficient has lower bound eight from the stronger
filtration. A total-order-three correction loses at most one bit,
so it stays above the leading valuation six, while later orders gain
still more. The same argument proves exactness for both mixed steps.
\end{proof}

\begin{remark}[What the power orders do not decide]
The exact orders in Theorem~\ref{thm:two-exact-power-orders} form an
infinite family. They do not supply a complete valuation formula at
non-power-of-two orders in address two, nor an exact higher-difference
formula at address three. Numerical patterns at other orders remain
bounded observations unless proved separately.
\end{remark}

\section{Exact first differences, least periods, and dyadic images}
\label{sec:sharp-periods-final}

The factorial estimates are stronger than needed for the optimal
first-difference law. A general Newton criterion gives the exact law
from the threshold-free bound and finitely many rational coefficients.
This route also yields all addresses at ranks two and eight, and an
infinite family at every larger dyadic rank.

\begin{lemma}[A first-difference criterion]
\label{lem:newton-exact-isometry-criterion}
Let $C:\mathbb Z_{\geq0}\to\mathbb Z_2$ have Newton coefficients
$a_k=\Delta_1^kC(0)$. Suppose for an integer $b\geq0$ that
\begin{equation}
 v_2(a_1)=b,\qquad
 v_2(a_k)\geq b+1+\lfloor\log_2k\rfloor\quad(k\geq2).
 \label{eq:newton-isometry-criterion}
\end{equation}
Then for every $n\geq0$ and $H>0$,
\begin{equation}
 v_2(C(n+H)-C(n))=b+v_2(H).
 \label{eq:newton-isometry-conclusion}
\end{equation}
The Newton interpolation extends continuously to $\mathbb Z_2$ and
satisfies the same identity for any two distinct dyadic arguments.
\end{lemma}

\begin{proof}
At a natural argument the exact Newton expansion is finite.
Vandermonde gives
\[
 \binom{n+H}{k}-\binom nk
   =\sum_{i=1}^k\binom Hi\binom n{k-i}.
\]
For $1\leq i\leq H$,
$v_2(\binom Hi)\geq v_2(H)-v_2(i)
 \geq v_2(H)-\lfloor\log_2k\rfloor$. Terms with $i>H$ vanish.
Thus each Newton summand with $k\geq2$ has valuation at least
$b+v_2(H)+1$. The $k=1$ summand is $Ha_1$, of exact valuation
$b+v_2(H)$, proving \eqref{eq:newton-isometry-conclusion}.

The coefficient hypothesis implies $a_k\to0$ dyadically, so the
Newton series converges uniformly on $\mathbb Z_2$, where
$\binom xk$ is integral-valued. For distinct $x,y\in\mathbb Z_2$,
approximate the pair by natural integer pairs with their difference
valuation eventually equal to $v_2(x-y)$. The corresponding function
differences have the fixed exact valuation $b+v_2(x-y)$.
A limit of elements with that fixed valuation has the same valuation:
the normalized residues are units, and $\mathbb Z_2^\times$ is closed.
\end{proof}

\begin{theorem}[Exact four-address first differences]
\label{thm:rank-four-exact-periods}
For every $n\geq0$, $H>0$, and $j\in\{0,1,2,3\}$,
\begin{align}
 v_2(C_j(n+H)-C_j(n))&=v_2(H)+b_j,
      &(b_0,b_1,b_2,b_3)&=(1,1,4,2),\label{eq:rank-four-normalized-isometry}\\
 v_2(E_{4(n+H)}^{(4,j)}-E_{4n}^{(4,j)})&=v_2(H)+B_j,
      &(B_0,B_1,B_2,B_3)&=(1,1,5,3).
 \label{eq:rank-four-original-isometry}
\end{align}
Consequently, the least positive row period modulo $2^r$ is exactly
\begin{equation}
 P_j(r)=2^{\max(0,r-B_j)}\qquad(r\geq1).
 \label{eq:rank-four-least-period-final}
\end{equation}
The corresponding displacement in the original Euler subscript is
$4P_j(r)$.
\end{theorem}

\begin{proof}
The actual reciprocal recurrence is
\begin{equation}
 C_{N,j}(n)=-\sum_{q=1}^n\binom{Nn}{Nq}
           g_{N,j}(q)C_{N,j}(n-q),\qquad C_{N,j}(0)=1.
 \label{eq:exact-reciprocal-row-recurrence}
\end{equation}
For $N=4$, the four kernels are
\[
 1,\qquad\frac1{4q+1},\qquad
 \frac1{(4q+1)(2q+1)},\qquad
 \frac3{(4q+1)(2q+1)(4q+3)}.
\]
The threshold-free theorem gives $v_2(a_k)\geq k$ in every address.
The following exact values cover all orders where that bound does not
already imply \eqref{eq:newton-isometry-criterion}:
\begin{center}
\begin{tabular}{c r l r}
$j$&$k$&$a_k$&$v_2(a_k)$\\\hline
0&1&$-2$&1\\
0&2&$72$&3\\
1&1&$-6/5$&1\\
1&2&$184/45$&3\\
2&1&$-16/15$&4\\
2&2&$64/45$&6\\
2&3&$-14848/1365$&9\\
2&4&$71426048/69615$&13\\
2&5&$-25509691392/85085$&16\\
2&6&$149113198870528/765765$&22\\
3&1&$-36/35$&2\\
3&2&$256/231$&8\\
3&3&$-65792/35035$&8\\
3&4&$161644544/4849845$&15
\end{tabular}
\end{center}
They are finite rational evaluations of
\eqref{eq:exact-reciprocal-row-recurrence}, followed by
$a_k=\sum_{h=0}^k(-1)^{k-h}\binom kh C_j(h)$.
For addresses zero and one, $k\geq3$ is covered by the generic bound.
For address two, $k\geq7$ is covered. For address three, $k\geq5$ is
covered. In each range the elementary inequality is
$k-\lfloor\log_2k\rfloor\geq b_j+1$.
Lemma~\ref{lem:newton-exact-isometry-criterion} proves the normalized
identity. Multiplication by $j!$ proves the original identity.

A row displacement $H$ preserves the original coefficient modulo
$2^r$ at even one row exactly when $v_2(H)+B_j\geq r$.
The least positive such integer is
$2^{\max(0,r-B_j)}$, and the same displacement works at every row.
This proves both optimality and the absence of a starting threshold.
\end{proof}

\begin{corollary}[Exact image and residue enumeration]
\label{cor:four-exact-images}
Let $\widetilde C_j$ denote the dyadic Newton interpolation.
The map
\[
 T_j(x)=\frac{\widetilde C_j(x)-1}{2^{b_j}}
\]
is a bijective isometry $\mathbb Z_2\to\mathbb Z_2$, and
\begin{equation}
 \widetilde C_j(\mathbb Z_2)=1+2^{b_j}\mathbb Z_2,
 \qquad j!\widetilde C_j(\mathbb Z_2)=j!+2^{B_j}\mathbb Z_2.
 \label{eq:four-interpolation-images-final}
\end{equation}
For $r>B_j$, one least period modulo $2^r$ visits every residue in
$j!+2^{B_j}\mathbb Z/2^r\mathbb Z$ exactly once. For $r\leq B_j$,
there is one residue and the least row period is one.
\end{corollary}

\begin{proof}
The exact difference identity gives
$v_2(T_j(x)-T_j(y))=v_2(x-y)$ and $T_j(0)=0$.
For every $s\geq1$, the map induces an injection of
$\mathbb Z_2/2^s\mathbb Z_2$ into itself. It is a bijection because
both sets have $2^s$ elements. These quotient bijections are compatible
with reduction. Inverting them gives a compatible preimage system for
every dyadic integer, proving surjectivity. Multiplication by $j!$
changes the image ideal by exactly $v_2(j!)$.
The finite quotient bijection is also precisely the residue enumeration
in a least row period.
\end{proof}

\subsection{An infinite exact family at all dyadic ranks}

\begin{theorem}[Lower-half exact laws]
\label{thm:lower-half-all-dyadic-isometry}
Let $N=2^a\geq4$ and $0\leq j<N/2$. Then for all $n\geq0$ and $H>0$,
\begin{equation}
 v_2(C_{N,j}(n+H)-C_{N,j}(n))=1+v_2(H).
 \label{eq:lower-half-all-dyadic-exact}
\end{equation}
Its normalized interpolation is a bijection onto $1+2\mathbb Z_2$.
For the original coefficients, the least row period modulo $2^r$ is
$2^{\max(0,r-1-v_2(j!))}$.
\end{theorem}

\begin{proof}
Every slope $N/i$ in the source has valuation at least two, since
$i\leq j<N/2$. Hence $g(1)\equiv1\pmod4$ and
$g(2)\equiv1\pmod8$. The first two Newton coefficients from the
reciprocal recurrence are
\begin{equation}
 a_1=-g(1)-1,\qquad
 a_2=\binom{2N}{N}g(1)^2-g(2)+2g(1)+1.
 \label{eq:lower-half-initial-newton}
\end{equation}
The first has valuation exactly one. Also
$\binom{2N}{N}\equiv6\pmod8$ for every $N=2^a\geq2$.
To see this, write
\[
 B_a=\binom{2^{a+1}}{2^a}
       =2\prod_{h=1}^{2^a-1}\left(1+\frac{2^a}{h}\right).
\]
The even-$h$ factors give $B_{a-1}$. The odd-$h$ factors are one
modulo four for $a\geq2$. Since $B_{a-1}$ has valuation one,
$B_a\equiv B_{a-1}\pmod8$, and $B_1=6$.
Therefore $a_2\equiv6-1+2+1=0\pmod8$.
The global Mahler bound covers all $k\geq3$, completing the criterion
with $b=1$. The period and image deductions are the same as above.
The proof also covers $N=2,j=0$, whose kernel is identically one.
\end{proof}

\subsection{Complete exact laws at ranks two and eight}

\begin{proposition}[Small-rank complete isometry table]
\label{prop:two-four-eight-isometry-table}
Every address at $N\in\{2,4,8\}$ satisfies
\[
 v_2(C_{N,j}(n+H)-C_{N,j}(n))=b_{N,j}+v_2(H).
\]
The following entries are ordered by $j=0,\ldots,N-1$:
\begin{center}
\begin{tabular}{c l l}
$N$&normalized shifts $b_{N,j}$&original shifts $B_{N,j}=b_{N,j}+v_2(j!)$\\\hline
2&$1,2$&$1,2$\\
4&$1,1,4,2$&$1,1,5,3$\\
8&$1,1,1,1,4,3,2,2$&$1,1,2,2,7,6,6,6$
\end{tabular}
\end{center}
For each entry, the least original row period is
$2^{\max(0,r-B_{N,j})}$, its Euler subscript displacement is $N$ times
that period, and the exact-image and residue-enumeration conclusions
hold with the corresponding shifts.
\end{proposition}

\begin{proof}
Rank four was proved above. At ranks two and eight the finite
criterion is verified by the actual rational recurrence
\eqref{eq:exact-reciprocal-row-recurrence}. The following valuation
lists record every initial Newton coefficient needed in addition to
the generic tail $v_2(a_k)\geq k$:
\begin{center}
\begin{tabular}{c c c l}
$N$&$j$&$b_{N,j}$&$v_2(a_1),v_2(a_2),\ldots$ through the finite cutoff\\\hline
2&0&1&$1,3$\\
2&1&2&$2,5,9,11$\\
8&0&1&$1,3$\\
8&1&1&$1,3$\\
8&2&1&$1,3$\\
8&3&1&$1,3$\\
8&4&4&$4,6,9,13,16,20$\\
8&5&3&$3,6,10,13,17$\\
8&6&2&$2,6,8,13$\\
8&7&2&$2,11,8,15$
\end{tabular}
\end{center}
The first entries can also be read directly from
$a_1=-1-\binom{N+j}{j}^{-1}$. At rank eight they are respectively
\[
 -2,\ -\frac{10}{9},\ -\frac{46}{45},\ -\frac{166}{165},\
 -\frac{496}{495},\ -\frac{1288}{1287},\
 -\frac{3004}{3003},\ -\frac{6436}{6435}.
\]
All later entries in the table are finite rational arithmetic, with
full numerator and denominator records retained in the accompanying
exact certificate \texttt{all\_rank\_extensions\_verified.json}.
Every entry with $k\geq2$ meets
$b_{N,j}+1+\lfloor\log_2k\rfloor$. Beyond the displayed cutoff, the
inequality $k\geq b_{N,j}+1+\lfloor\log_2k\rfloor$ holds at every
order. Thus the finite table and the proved infinite tail establish
the full hypothesis of Lemma~\ref{lem:newton-exact-isometry-criterion}.
This is a finite coefficient proof followed by a universal tail
argument, rather than extrapolation from sampled rows or steps.
\end{proof}

\begin{remark}[Failure of a universal upper-half exact law]
The all-rank lower bound remains valid in every address. The exact
first-difference law does not automatically extend to all upper-half
addresses. Exact rational recurrence calculations give
\begin{center}
\begin{tabular}{c c c c}
$(N,j)$&$v_2(C(1)-C(0))$&$v_2(C(2)-C(0))$&naive $b+v_2(2)$\\\hline
$(16,9)$&5&7&6\\
$(32,18)$&6&6&7\\
$(32,19)$&5&7&6
\end{tabular}
\end{center}
For example, at $(N,j)=(16,9)$ the first difference is
$-2042976/2042975$, whereas the two-row difference is
$-2215603929728/2215923048625$. These exact counterexamples refute
the additional universal isometry proposal. They preserve the
threshold-free lower bound and the original conjecture.
\end{remark}

\section{Analytic interpolation and exact Taylor radii}
\label{sec:analytic-radii-final}

The threshold-free Mahler theorem gives continuous interpolation.
Analytic interpolation needs a stronger estimate, and the factorial
bounds provide it directly. The ordinary Taylor radius is determined
by the coefficient growth, including infinite exact subsequences.

\begin{theorem}[Taylor coefficient bounds and radii]
\label{thm:four-taylor-radii}
For $j=0,1,2,3$, the dyadic interpolation has an ordinary power series
\[
 \widetilde C_j(x)=1+\sum_{\ell\geq1}d_{j,\ell}x^\ell,
 \qquad v_2(d_{j,\ell})\geq\ell+\beta_j.
\]
Its Taylor series converges for $|x|_2<2$.
At addresses zero and one,
\begin{equation}
 v_2(d_{j,\ell})=\ell\quad(\ell\geq1),
 \qquad R_0=R_1=2.
 \label{eq:zero-one-exact-taylor-radius}
\end{equation}
At address two,
\begin{equation}
 v_2(d_{2,\ell})\geq
       \max\left(2\ell+\left\lfloor\frac\ell2\right\rfloor,
                          \ell+3\right),
 \qquad R_2=2^{5/2}.
 \label{eq:two-exact-taylor-radius}
\end{equation}
The radius at address three is at least two. Its exact radius is not
asserted here. The stated exact radii at addresses zero, one, and two
hold at every center in $\mathbb Z_2$.
\end{theorem}

\begin{proof}
Put $a_m=\Delta_1^mC_j(0)$ and expand the falling factorial in the
Newton series. The ordinary coefficients are
\begin{equation}
 d_{j,\ell}=\sum_{m\geq\ell}\frac{a_m}{m!}s(m,\ell).
 \label{eq:newton-taylor-stirling}
\end{equation}
Since $s(m,\ell)$ is integral and
$v_2(a_m/m!)\geq m+\beta_j$, the coefficient sum converges and has
valuation at least $\ell+\beta_j$. The sum exchange is justified on
$\mathbb Z_2$ by these uniform geometric coefficient bounds.
The ordinary power series agrees there with the Newton interpolation
and converges whenever $|x|_2<2$.

For $j=0,1$, the term $m=\ell$ in
\eqref{eq:newton-taylor-stirling} has exact valuation $\ell$, since
$s(\ell,\ell)=1$ and the Newton valuation is exact. Every later term
has valuation at least $\ell+1$. Thus
$v_2(d_{j,\ell})=\ell$, and the coefficient radius is exactly two.
At a center $x_0\in\mathbb Z_2$, the coefficient of $y^\ell$ in the
translated series is
\begin{equation}
 \sum_{m\geq\ell}d_{j,m}\binom m\ell x_0^{m-\ell}.
 \label{eq:taylor-translation-coefficient}
\end{equation}
Its initial term again has exact valuation $\ell$, and all later
terms have at least one more bit. The exact radius is therefore two
at every integral center.

For address two, the stronger factorial filtration gives the lower
bound in \eqref{eq:two-exact-taylor-radius}. At each
$\ell=2^r\geq4$, Theorem~\ref{thm:two-exact-power-orders} makes the
$m=\ell$ term in \eqref{eq:newton-taylor-stirling} have valuation
$2\ell+\ell/2=5\ell/2$. Every $m>\ell$ term has strictly greater
lower bound $2m+\lfloor m/2\rfloor$. Hence
\[
 v_2(d_{2,\ell})=5\ell/2\qquad(\ell=2^r\geq4).
\]
The universal lower bound and this unbounded exact subsequence prove
that the coefficient radius is exactly $2^{5/2}$. Translation by an
integral center preserves the same unique leading coefficients at
these degrees in \eqref{eq:taylor-translation-coefficient}, proving
the exact radius at every such center.
\end{proof}

\begin{remark}[Boundary interpretation]
For addresses zero and one, at an argument of absolute value two every
positive-degree Taylor term has valuation zero, so that ordinary Taylor
series diverges. For address two, over a complete extension admitting
arguments of absolute value $2^{5/2}$, the power-of-two-degree terms
have valuation zero and fail to tend to zero. These statements concern
the Taylor expansions just constructed. They do not prohibit a
different analytic continuation. On $\mathbb Q_2$-valued arguments,
open-disc notation must be read with its discrete value group. The
coefficient radius is the precise invariant in the theorem.
\end{remark}

\section{The final congruence in the original Euler convention}
\label{sec:final-original-congruence}

The preceding bounds can be stated as one uniform congruence without
suppressing the address normalization. This is the strengthened formula
for the four-address family, valid from the first row.

\begin{corollary}[Final common-step congruence]
\label{cor:final-euler-congruence}
Let $j\in\{0,1,2,3\}$, $n\geq0$, and $k,H\geq1$. Define
\[
 A_0(k)=A_1(k)=k,\qquad
 A_2(k)=\max\left(k+4,\ 2k+\left\lfloor\frac k2\right\rfloor+1\right),
 \qquad A_3(k)=k+2.
\]
Then
\begin{equation}
 \boxed{\displaystyle
 \sum_{i=0}^k(-1)^{k-i}\binom ki E_{4(n+iH)}^{(4,j)}
       \equiv0\pmod{2^{\,v_2(k!)+k v_2(H)+A_j(k)}}.}
 \label{eq:boxed-final-euler-congruence}
\end{equation}
For $j=0,1$, this displayed exponent is the exact valuation for every
$n,k,H$ in the stated range. For $j=2$, the exact positive
power-of-two-order valuations are given by
Theorem~\ref{thm:two-exact-power-orders}, with one added bit.
\end{corollary}

\begin{proof}
The left side is $j!\Delta_H^kC_j(n)$.
At addresses zero and one, $v_2(j!)=0$ and the exact homogeneous
formula gives the exponent. At address two, $v_2(2!)=1$. The maximum
of \eqref{eq:factorial-four-lower} and
\eqref{eq:factorial-four-j2-filtration} yields
$\max(k+4,2k+\lfloor k/2\rfloor+1)$ after removing the common
$v_2(k!)+k v_2(H)$. At address three, the one-bit normalized offset
and $v_2(3!)=1$ give $k+2$.
\end{proof}

For mixed steps the same corollary holds with $k v_2(H)$ replaced by
$\sum_\nu v_2(H_\nu)$ and the binomial stencil replaced by
$j!\Delta_{H_1}\cdots\Delta_{H_k}C_j(n)$.
At addresses zero and one, exactness is then governed by the parity
classification in Theorem~\ref{thm:zero-one-exact-mixed}.
At first order these bounds are exact in every address, with original
valuation offsets $(1,1,5,3)$. The least periods are those in
\eqref{eq:rank-four-least-period-final}.

\subsection{Proof coverage and remaining questions}

The results in these sections are mathematical theorems with detailed
proofs and preserved audit evidence. Their formal status must be read
statement by statement. The recovered Paper~IX release already contains
Lean results for the rank-four threshold-free homogeneous bound, the
sharp first-difference laws, the least periods, and the dyadic isometry
and image conclusions. Those inherited results should not be described
as missing merely because the earlier short completion supplement did
not include their later modules.

The all-rank dyadic and general prime-power transfer proof presented
here, the factorial improvements, the mixed leading-parity
classification, the Bernoulli power-order sharpness, the all-rank
lower-half exact laws, the complete rank-eight extension, and the exact
Taylor-radius arguments have not received new Lean formalization in
this strengthening. A mathematical audit and an exact computation are
neither a Lean compilation nor external peer review. The separate Lean
guide records the recovered formal declarations and any fresh release
checks without extending their coverage to these new claims.

The completed external statement is Nishibuchi's Conjecture~3.9 in
its prescribed parameter family and coefficient normalization.
These stronger congruences enlarge its dyadic branch and remove its
starting threshold there. They do not settle every Euler problem.
In particular, the following questions remain separate:
\begin{itemize}
 \item a complete exact higher-difference law at all non-power-of-two
       orders in address two;
 \item complete higher-difference valuations and the exact ordinary
       Taylor radius in address three;
 \item a classification of the upper-half dyadic addresses having an
       exact first-difference metric law at arbitrary larger ranks;
 \item the distinct universal simplicity or nonvanishing question
       for the later Euler-tail residuals, and the original native
       Paper~IX conjectures beyond the completed external statement.
\end{itemize}
The exact counterexamples above rule out the naive universal upper-half
isometry, and mixed cancellations rule out unconditional baseline
sharpness for arbitrary stencils. No global historical-priority or
novelty assertion is made for these strengthened statements.
